\chapter{Conclusión}

% ===========================================================================
% BOSQUEJO DE TRABAJO — Por desarrollar tras la validación del Capítulo 5.
%
% La mayoría de las secciones puede escribirse a partir de lo ya documentado
% (objetivos del Capítulo 1, contribuciones implícitas en Capítulos 3 y 4,
% limitaciones reconocidas en Capítulos 3 y 4). La sección 6.5 es opcional
% y depende de la reflexión personal una vez cerrado el resto.
% ===========================================================================

% ---------------------------------------------------------------------------
% 6.1 Síntesis del trabajo realizado
% ---------------------------------------------------------------------------
% - Resumir qué se construyó: plataforma educativa autohospedable de
%   ciberseguridad ofensiva con validación automática basada en estado del
%   sistema, motor de evaluación con targets atómicos y nodos lógicos
%   (and/or/not), y daemon REPL multi-cliente para uso educativo.
% - Recorrer los OBJETIVOS ESPECÍFICOS del Capítulo 1 y marcar el estado:
%     1. Arquitectura modular multi-entorno con virtualización   (cumplido).
%     2. Configuración declarativa YAML accesible al educador    (cumplido).
%     3. Motor de validación automática basado en estado         (cumplido).
%     4. Herramientas CLI de seguimiento y administración        (cumplido).
%     5. Poner la plataforma a disposición del CLCERT            (parcial:
%        depende de la entrega y las pruebas en su infraestructura — ver
%        Capítulo 5).

% ---------------------------------------------------------------------------
% 6.2 Contribuciones
% ---------------------------------------------------------------------------
% Tres ejes a desarrollar:
%
% - PEDAGÓGICO: trazabilidad real del progreso del estudiante. En lugar de un
%   flag binario "completó / no completó", el educador puede ver QUÉ acción
%   concreta disparó cada target gracias al trigger_event persistido en cada
%   TargetResult, y reconstruir la secuencia del ataque. Esto cambia la
%   pregunta de "¿lo logró?" a "¿cómo lo logró?".
%
% - TÉCNICO Y DE COMUNIDAD: contribución al ecosistema open-source de
%   detección de intrusiones. El código, los patrones de detección
%   (network_payload, process_running con ancestros, sidecar tcpdump por
%   escenario, etc.) y la metodología están en un repositorio público y
%   auditable, frente a las reglas cerradas de los IDS comerciales.
%
% - ARQUITECTÓNICO: separación clara en 4 capas con contratos explícitos.
%   Esto hace al sistema extensible (nuevos tipos de target en Capa 3 sin
%   tocar Capa 1) y portable (cambio de tecnología de virtualización
%   implementando solo EnvironmentProvider). El precedente más cercano,
%   Tectonic, mezcla estas responsabilidades; FlexiPwn las separa.

% ---------------------------------------------------------------------------
% 6.3 Limitaciones reconocidas
% ---------------------------------------------------------------------------
% Buena parte ya está documentada en los capítulos previos y se puede citar
% sin reabrir la discusión:
%
% - GAP DEL POLLING: procesos sub-segundo, fileless, PPID spoofing y técnicas
%   evasivas avanzadas quedan fuera del alcance de detección. Aceptable para
%   el contexto educativo. (Detalle en Capítulo de Implementación, Sección
%   "Limitaciones del Monitoreo por Polling").
%
% - LIMITACIONES DEL NETWORKMONITOR:
%     * Tráfico cifrado (TLS) y protocolos binarios solo exponen metadatos.
%     * Fragmentación de payload entre paquetes TCP — el parser opera por
%       paquete completo, sin reensamblado a nivel de flujo.
%
% - LIMITACIÓN DEL LOGMONITOR: dependencia del flush de la aplicación;
%   exige cooperación del Dockerfile del educador (PYTHONUNBUFFERED=1 o
%   equivalente).
%
% - DEUDA TÉCNICA HEREDADA DE ITERACIONES TEMPRANAS, que queda como
%   limpieza pendiente en el repositorio:
%     * Tipos de target file_exists, http_response_contains y
%       database_query_result definidos en el schema Pydantic pero sin
%       evaluador productor — se eliminarán del schema (no se implementarán).
%     * Campo attacker_ports del EnvironmentConfig sin uso real: el binding
%       del SSH del atacante se asigna automáticamente desde el código.
%     * Dependencia watchdog declarada en pyproject.toml sin uso en código
%       (residuo de la implementación inicial con inotify, sustituida por
%       container.diff()).
%
% - VALIDACIÓN EN INFRAESTRUCTURA REAL: pendiente al momento de redactar.
%   La eficacia del sistema en aula multi-estudiante y la robustez bajo carga
%   están parcialmente cubiertas — ver Capítulo 5.

% ---------------------------------------------------------------------------
% 6.4 Trabajo futuro
% ---------------------------------------------------------------------------
% Organizado en dos horizontes temporales:
%
% CORTO PLAZO — Limpieza y correcciones ya identificadas
%
% - Remover los tipos de target sin productor (file_exists,
%   http_response_contains, database_query_result) del schema Pydantic y de
%   todo material derivado.
% - Remover el campo attacker_ports del EnvironmentConfig.
% - Remover la dependencia watchdog del pyproject.toml.
% - Rehacer el diagrama TikZ de arquitectura para reflejar la interfaz
%   abstracta + tecnología concreta en Capa 1 y agregar el monitor de red en
%   Capa 2.
% - Cerrar la entrega al CLCERT con documentación operativa.
%
% MEDIANO Y LARGO PLAZO — Extensiones del sistema
%
% - PodmanProvider como segundo backend de la interfaz EnvironmentProvider
%   (justificación detallada en Capítulo de Implementación, sección de
%   virtualización). Mejora arquitectónica de seguridad sin afectar capas
%   superiores.
% - INTERFAZ WEB que exponga las mismas operaciones que la CLI clásica
%   manteniendo la lógica de negocio intacta. Útil especialmente para el
%   CLCERT y estudiantes con menos manejo de línea de comandos.
% - DETECCIÓN BASADA EN eBPF / Falco / Sigma rules como complemento del
%   polling, para cerrar el gap de detección de procesos sub-segundo y
%   técnicas evasivas (fileless, PPID spoofing).
% - KATA CONTAINERS o gVisor como runtime alternativo para ejercicios de
%   kernel exploitation, donde el aislamiento de Docker rootless no alcanza.
% - REENSAMBLADO TCP en NetworkMonitor para soportar payloads fragmentados
%   entre paquetes (necesario para escenarios con tráfico de mayor volumen).
% - DOCUMENTACIÓN para creadores de escenarios + plantilla de Dockerfile
%   recomendada, cubriendo el caso del flush y otras buenas prácticas que
%   evitan los puntos ciegos del LogMonitor.

% ---------------------------------------------------------------------------
% 6.5 Reflexión final  (una página corta)
% ---------------------------------------------------------------------------
% - Lugar del proyecto en el ecosistema open-source de cyber ranges
%   educativos en español, y vínculo con el contexto de la universidad y el
%   CLCERT.
% - Aprendizajes personales: lo más valioso de construirlo, lo que se haría
%   distinto si se empezara otra vez.
